\subsection{Impact of the Number of Switches per ISP}
\label{subsec:isp}

In the first experiment, the number of switches per ISP is increased from 5 to 10. The requests are limited to 6--8 virtual nodes and a bandwidth demand of 10 or 15. Throughout this section, the acceptance ratio is computed over all requests, whereas the network resource usage, the average path length and the end-to-end delay are averaged over the common requests, i.e., the requests accepted by all four algorithms, so that all algorithms are compared on identical requests. Figs.~\ref{fig:isp_nsf_acc} and \ref{fig:isp_us_acc} show that the main effect of the number of switches appears in the acceptance ratio. With 5 switches, the greedy approaches accept only about 60\% of the requests in NSFNET and 75--85\% in USNET. In contrast, the GA-based algorithms exceed 97\% in both networks. When the number of switches increases to 6, the number of pathlets offered by the domains in NSFNET grows by about 3.8 times (2.3 times in USNET). Accordingly, the acceptance ratio rises to 87.5\% for CRG-VMTE and 93.3\% for CAN-VMTE. However, this improvement is not monotonic. In USNET, the acceptance ratio of CRG-VMTE drops from 95.8\% to 80.8\% when the number of switches increases from 7 to 8. In the same step, CAN-VMTE drops from 98.3\% to 88.3\%, while QAG-VMTE remains at 100\%. CAN-VMTE shows a similar drop in USNET from 5 to 6 switches (from 85.0\% to 72.5\%), and both greedy approaches decline slightly in NSFNET from 6 to 7 switches. Since the number of pathlets in USNET still grows by 13\% from 7 to 8 switches, the drop is not caused by fewer pathlets. This drop is likely caused by bottlenecks in the topology instances generated for this scale. The GA-based algorithms are not affected, which indicates that a wider search space improves robustness.

Figs.~\ref{fig:isp_nsf_res} and \ref{fig:isp_us_res} show that QAG-VMTE achieves the lowest resource usage at all scales. Compared with CRG-VMTE, it uses 45--57\% fewer resources in NSFNET and 47--56\% fewer resources in USNET. Although more alternative paths become available with more switches, the path length and delay do not decrease significantly. For the greedy approaches, they even increase; for example, the path length of CRG-VMTE in NSFNET grows from 38.1 at 6 switches to 41.6 at 10 switches, whereas that of QAG-VMTE stays between 16.0 and 18.0. Compared with CRG-VMTE, QAG-VMTE reduces the delay by 34--44\% in USNET. In NSFNET, it reduces the delay by 36--40\% from 6 switches onward. With 5 switches in NSFNET, this advantage drops to 29.8\%, because the common set then contains only the 314 requests that both greedy approaches can embed. When more requests become common at 6 switches (501 requests), the delay of CRG-VMTE on the common set increases by 13.9\%, whereas that of QAG-VMTE slightly decreases from 185.3 to 182.0. Hence, the additional requests are harder for CRG-VMTE but not for QAG-VMTE.

\subsection{Impact of the Bandwidth Demand}
\label{subsec:vbw}

In the second experiment, the bandwidth demand is increased from 5 to 25 and the results of the 6--8 switch scales are combined. Figs.~\ref{fig:vbw_nsf_acc} and \ref{fig:vbw_us_acc} show that all algorithms accept more than 96\% of the requests when the demand is 10 or lower, and they start to diverge at 15. At a demand of 25, the acceptance ratio of CRG-VMTE falls to 31.3\% in NSFNET and 27.5\% in USNET. QAG-VMTE stays above 97\% in both networks and reaches 3.2 to 3.6 times the acceptance ratio of CRG-VMTE. The gap between GARC-VMTE and QAG-VMTE also peaks at this point and reaches 7.7 percentage points in NSFNET and 8.1 percentage points in USNET. This result shows that the contribution of Q-learning becomes clearer as the capacity pressure increases.

As shown in Figs.~\ref{fig:vbw_nsf_res} and \ref{fig:vbw_us_res}, the network resource usage increases almost linearly with the demand by definition. However, the advantage of QAG-VMTE grows as the demand increases. In NSFNET, its advantage over CRG-VMTE increases from 52\% to 56\% and its advantage over GARC-VMTE increases from 25\% to 32\%. In USNET, these advantages increase from 51\% to 55\% and from 21\% to 30\%, respectively. In USNET, the slower increase at a demand of 25 is not a capacity gain. From a demand of 20 to 25, the resource usage grows by only 3--5\% for the greedy approaches and QAG-VMTE (14\% for GARC-VMTE), although the demand itself grows by 25\%. It is a sampling effect caused by the drop in the number of common requests from 118 to 55, which also reduces the mean number of virtual nodes of the common requests from 6.7 to 6.0. In NSFNET, the common set shrinks less (from 773 to 393 requests, with the mean number of virtual nodes decreasing from 7.1 to 6.5), and the resource usage still grows by 20--25\%. The same composition effect shapes the path length and the delay. Over the common requests, both metrics remain almost constant up to a demand of 15 and then decrease for all algorithms, by 4--9\% in NSFNET and 17--21\% in USNET from a demand of 15 to 25. The relative differences, however, remain stable: at every demand level, QAG-VMTE reduces the delay by 36--39\% in NSFNET and by 38--41\% in USNET compared with CRG-VMTE.

\subsection{Impact of the Number of Virtual Nodes}
\label{subsec:vnode}

In the third experiment, the number of virtual nodes per request is increased from 5 to 10 and the results of the 6--8 switch scales are combined. Figs.~\ref{fig:vnode_nsf_acc} and \ref{fig:vnode_us_acc} show that the acceptance ratio of the greedy approaches decreases gradually rather than sharply. With 10 virtual nodes, the acceptance ratio of CRG-VMTE falls to 62.9\% in NSFNET and 58.3\% in USNET. The acceptance ratio of QAG-VMTE does not fall below 98.3\% in either network. This effect becomes stronger as the bandwidth demand increases. In NSFNET at a demand of 25, the acceptance ratio of CRG-VMTE drops from 50\% with 5 virtual nodes to 14\% with 10 virtual nodes. In the same case, QAG-VMTE maintains about 98\%. In USNET, the corresponding drop of CRG-VMTE is from 53\% to 5\%, while QAG-VMTE stays between 95\% and 100\%.

Figs.~\ref{fig:vnode_nsf_len}--\ref{fig:vnode_nsf_delay} and \ref{fig:vnode_us_len}--\ref{fig:vnode_us_delay} show that the number of virtual nodes is the most influential parameter on path length and delay. Since the number of virtual links increases from 4 to 9, the path length more than doubles for all algorithms. However, the number of hops per virtual link remains almost constant. In NSFNET, it is about 2.7 for QAG-VMTE and 3.8 for GARC-VMTE. For the greedy approaches, it lies between 5.4 and 6.0. In USNET, these values are about 4.0, 5.2 and 7.9--8.5, respectively. Therefore, all algorithms preserve their path quality as the requests grow. The fastest relative increase of QAG-VMTE is due to its very low initial value. Delay increases more slowly than path length (by a factor of 1.5--1.7 from 5 to 10 virtual nodes, compared with 2.1--2.4 for path length). This is because path length is summed over all virtual links, whereas delay is determined only by the longest end-to-end path. Across all virtual node counts, QAG-VMTE uses 53--57\% fewer resources than CRG-VMTE in NSFNET and 50--55\% fewer resources in USNET. It also reduces the delay by 37--41\% and 37--44\% in these networks, respectively.

\subsection{Overall Evaluation and the Effect of Q-Learning}
\label{subsec:overall}

Over the entire parameter space, the acceptance ratio in NSFNET is 73.5\% for CRG-VMTE, 78.3\% for CAN-VMTE, 96.2\% for GARC-VMTE and 98.5\% for QAG-VMTE. In USNET, these values are 74.3\%, 77.0\%, 96.7\% and 98.9\%, respectively. The proposed algorithms also make fewer inter-domain transitions than the greedy approaches. From CRG-VMTE to GARC-VMTE, the number of inter-domain transitions decreases by 37--40\% and the delay decreases by 32--35\%. In the larger USNET, the number of hops per virtual link increases by 31--39\% for all algorithms (e.g., from 6.2 to 8.5 for CRG-VMTE and from 2.8 to 3.9 for QAG-VMTE), so the ranking of the algorithms is preserved. Among the greedy approaches, CAN-VMTE has no consistent advantage over CRG-VMTE. It uses fewer resources in only 58--62\% of the common requests, and it achieves a lower delay in 58\% of the common requests in NSFNET and in only 51\% in USNET, which is close to chance level.

The comparison between GARC-VMTE and QAG-VMTE directly reveals the effect of Q-learning. Compared with GARC-VMTE, Q-learning-based selection reduces the resource usage by 19--32\% in all experiments. It also uses fewer resources in more than 99.5\% of the common requests. Since the resource usage of a request is the product of its bandwidth demand and its hop count, and all algorithms are evaluated on the same requests, this reduction is almost identical to the reduction in path length. Its effect on the acceptance ratio appears only under high resource pressure, such as with few switches per domain, high bandwidth demand or large requests; the largest gains are 7.7 and 8.1 percentage points at a demand of 25 and 4.3 percentage points with 10 virtual nodes in USNET. The gain in delay is more limited. QAG-VMTE reduces the hop count by about 23--27\% but reduces the delay by only 7--9\%. It achieves a lower delay in 84\% of the common requests. The recorded paths explain this result. The two algorithms make almost the same number of inter-domain transitions (8.3 and 8.1 per request in NSFNET, 11.5 and 11.4 in USNET), and the gain of Q-learning comes mainly from intra-domain hops, which QAG-VMTE reduces by 35--40\%. This reduction largely stems from node placement. QAG-VMTE places 71\% (NSFNET) and 81\% (USNET) of the virtual nodes on edge routers, compared with 43--49\% for GARC-VMTE and 36--46\% for the greedy approaches, although edge routers constitute only 30--70\% of the physical nodes depending on the scale. When a virtual node is hosted on the edge router through which its virtual link leaves or enters the domain, no intra-domain hop is needed on that side. This holds for 57--58\% of the link ends of QAG-VMTE, but only for 11--19\% of those of the other algorithms. Since inter-domain links have much higher delay, this gain has a limited effect on the end-to-end delay. Hence, the end-to-end delay in multi-domain networks is largely determined by inter-domain transitions, and a delay-oriented improvement should primarily target them.
